\chapter{总结与展望}\label{chap:conclusion}

\section{总结}\label{sec:conclusion-summary}

本文围绕基于自然语言的机械臂多任务通用控制方法展开系统性研究，旨在探索轻量级语言模型作为机械臂通用控制器的可行性与技术路线。随着具身智能与自然语言处理技术的快速发展，以自然语言为交互接口实现机械臂的通用化控制已成为降低使用门槛、拓展应用场景的重要需求。在专家策略训练方面，本文基于TQC算法与HER机制训练了高性能的单任务专家策略，为后续数据集构建提供了高质量的轨迹来源。在数据集构建方面，本文设计了从物理状态到自然语言的异构数据转化机制与基于闭环回放的轨迹清洗方法，保障了数据集内轨迹的物理可复现性。在微调推理方面，本文在6个轻量级基座模型上开展了系统的LoRA微调实验，并借助vLLM推理引擎实现了低延迟闭环推理。仿真实验验证了所提方法在多任务控制中的有效性，实验结果表明基于LoRA微调的轻量LLM控制方案在多任务混合微调中取得了可接受的控制性能，跨任务知识迁移对多数模型具有正向增益，验证了单一LLM同时掌握多项操作技能的可行性。本文的研究成果为轻量大语言模型在机械臂控制领域的应用提供了实验参考与技术路线。

针对高质量专家轨迹数据的获取问题，本文提出了基于TQC算法与HER机制的机械臂单任务专家策略训练方案。该方案针对Fetch机械臂的Reach、Push、Slide与PickAndPlace四项操作任务，分别设计状态空间、动作空间与奖励函数，在MuJoCo仿真环境中训练独立的单任务专家策略。通过消融实验对比TQC与TQC+HER的训练表现，揭示了HER机制的双面性：在稀疏奖励环境中能有效加速训练收敛，但在目标空间分布与实际到达状态分布差异较大的任务中反而会降低策略网络最终性能。仿真验证了所训练专家策略在四项任务上的评估成功率分别达到100.0\%、97.9\%、99.0\%与90.6\%，为后续数据集构建提供了可靠的轨迹数据源。

针对物理状态与自然语言之间的跨模态数据转化问题，本文设计了从物理观测到结构化文本的异构数据转化机制与基于闭环回放的轨迹清洗方法。该机制通过对专家策略输出的数值观测轨迹进行语义切分与离散化，将其映射为具有物理含义的结构化字段，并通过三角色对话模板封装转化为语言模型可处理的文本格式。在此基础上，提出基于物理引擎闭环回放的轨迹采集与清洗方法，对转化后的数据进行可复现性验证与清洗。仿真验证表明，清洗后四项任务的轨迹保留率分别为100.0\%、90.6\%、92.6\%与70.0\%，任务复杂度与保留率呈负相关趋势，所有保留轨迹均通过物理引擎回放验证。

针对轻量级语言模型在多任务控制场景下的跨模态动作生成问题，本文利用构建的多任务专家数据集，通过LoRA参数高效微调在6个轻量级基座模型上开展了系统的微调实验。实验首先在单任务设定下评估各模型的动作预测精度，而后开展多任务混合微调实验，借助vLLM推理引擎实现闭环推理与环境测试。仿真结果表明，Qwen3.5-0.8B在多任务混合微调中以54.9\%的平均成功率取得最优性能。跨任务知识迁移的效果因模型而异：Gemma3-1B与Llama3.2-1B分别获得16.5与12.85个百分点的显著多任务增益，验证了单一LLM同时掌握多项操作技能的可行性；而Qwen3-0.6B出现了4.0个百分点的性能下降，表明简单的数据混合策略在模型容量有限时存在瓶颈。所有模型的推理延迟均在1\textasciitilde10 Hz范围内，验证了vLLM推理方案在实时推理层面的可行性。

\section{展望}\label{sec:conclusion-outlook}

本文工作在轻量LLM驱动的机械臂控制方向上取得了一定的实验进展，但仍存在局限，有待后续研究进一步完善。主要包括以下几点：

\begin{enumerate}
  \item 深入优化多任务微调策略。当前多任务混合微调以简单的数据混合方式实现，在参数量有限的条件下导致单项任务性能下降。后续工作可探索任务条件化的微调策略，如通过任务标识符前缀或路由机制引导模型在不同任务间动态切换参数子空间，以减少任务间的参数干扰。此外，扩大基座模型参数规模或增加LoRA秩也可能提升多任务编码容量。
  \item 融合视觉感知。当前系统的输入仅为结构化文本描述的物理状态，依赖仿真环境提供的精确状态观测，无法直接处理视觉输入。未来可引入视觉编码器，将图像观测与文本指令共同作为输入，使模型具备从原始视觉信号到动作的端到端映射能力，从而摆脱对精确状态估计的依赖，提升系统在真实场景中的适用性。
  \item 进行仿真到真实机器人的迁移。本文实验在MuJoCo仿真环境中完成，尚未验证在真实物理平台上的控制效果。仿真环境与真实世界之间存在动力学差异与传感器噪声，直接迁移可能面临性能衰减。后续可通过域随机化、系统辨识或少量真机数据二次微调等方法弥合仿真到真实的差距。
  \item 改进动作表示与生成机制。当前模型以自回归方式逐token生成动作向量的文本表示，生成过程与物理控制的连续性之间存在不匹配。引入扩散模型或流匹配等连续动作生成方法，可能提升动作输出的平滑性与多峰分布拟合能力，降低对输出格式的依赖。
  \item 扩展数据规模与任务多样性。当前数据集仅覆盖Fetch机械臂的四项基础操作任务，且每项任务的轨迹数量有限。扩大任务类型与数据规模有助于更全面地评估轻量LLM在机械臂控制中的泛化边界，也为研究数据质量对微调性能的影响规律提供更充分的实验基础。
\end{enumerate}
